% LSPV-1 REPORT::CUTGAP1::2026-10-12T14:14:00Z
\documentclass[11pt]{article}
\usepackage[margin=1in]{geometry}
\usepackage[T1]{fontenc}
\usepackage{amsmath,amssymb,amsthm,booktabs,longtable}
\usepackage{hyperref}
\newtheorem{theorem}{Theorem}
\newtheorem{lemma}{Lemma}
\newcommand{\OPEN}{\textbf{OPEN}}
\newcommand{\EXACT}{\textbf{PROVED/EXACT}}
\newcommand{\FALSIFIED}{\textbf{FALSIFIED/RETIRED}}
\newcommand{\KNOWN}{\textbf{KNOWN THEOREM}}
\newcommand{\COMP}{\textbf{COMPUTATIONAL}}
\begin{document}

\section*{CUTGAP1: polynomial scheduled cuts do not force compatibility exposure}

\paragraph{Executive status.}
\EXACT: a polynomial-size, uniformly predeclared schedule of normalized checkpoints can
be avoided by an extensionally equivalent polynomial wrapper.  The wrapper confines all
source-dependent work, including a simulation of the child decider and cleanup, to one
checkpoint-free interval.  At every scheduled cut its mutable state depends only on
encoded length and, after the hidden interval, the one-bit answer.  Hence no such state
vector is universally forced to determine a non-output semantic compatibility object.
This strictly extends the single-cut result of MODELBRIDGE-0026.  It is a negative
modeling theorem, not a superpolynomial lower bound: adaptive or exogenous cuts,
all-time/peak obligations, and a solver-independent transfer remain \OPEN.  P versus NP
remains open.

\section{Barrier-transfer literature preflight}

The two obstacles are (i) wrapper scheduling freedom and cleanup erase observations at
chosen internal times, and (ii) known SAT time--space lower bounds apply only on restricted
resource axes and do not force arbitrary polynomial-time deciders to materialize a chosen
semantic object.  Fortnow--van Melkebeek prove nontrivial SAT time--space tradeoffs on
random-access machines, not a superpolynomial lower bound against all polynomial time
\cite{fvm}.  Baker--Gill--Solovay exhibit relativized worlds on both sides of P versus NP
\cite{bgs}, and Aaronson--Wigderson show that even many arithmetizing methods face the
algebrization barrier \cite{aw}.  D\"ocker proves NP-completeness of the exact monotone
clean core used below \cite{docker}.

\begin{longtable}{p{.23\linewidth}p{.69\linewidth}}
\toprule
Transfer class & Assessment \\
\midrule
Directly transferable theorem & Exact-clock padding and deterministic cleanup are used
inside the same uniform deterministic machine model.  The elementary largest-gap lemma
below transfers without changing models. \\
Useful analogy requiring a new theorem & Cell-probe chronograms, distributed cut
arguments, control-theoretic observability, and VLSI clock/routing constraints become
relevant only when observation times or interfaces are externally fixed.  A theorem
making SAT computation obey such an exogenous interface is missing. \\
Misleading / failed transfer & Physical clocking, communication cuts, and proof-system
state lower bounds do not by themselves constrain an arbitrary single-tape or RAM SAT
decider.  Counting semantic states is not a running-time lower bound. \\
\bottomrule
\end{longtable}

\section{Explicit clean core and semantic compatibility}

Let $\mathcal K_N$ be the canonical encodings of duplicate-free monotone 3-CNF formulas
of encoded length $N$ in which every clause contains three distinct variables and each
variable occurs exactly twice positively and twice negatively.  Monotone means that each
clause is all-positive or all-negative.  SAT restricted to this family is NP-complete
\cite{docker}.  For a partial assignment $\alpha$ define
\[
 H(F,\alpha)=1 \quad\Longleftrightarrow\quad
 \exists z\supseteq\alpha\; F(z)=1.
\]
This is a semantic object, not a required internal representation.

The finite witness in the accompanying result file gives two six-variable, eight-clause
members $F_0,F_1\in\mathcal K_N$ with
\[
 \operatorname{SAT}(F_0)=\operatorname{SAT}(F_1)=1,
 \qquad H(F_0,x_1=x_2=0)=0,
 \qquad H(F_1,x_1=x_2=0)=1.
\]
Thus $H$ is nonconstant inside one equal-length YES fiber.  The complete clauses and
model sets are recorded, rather than inferred, in
\path{results/raw/review/2026-10-12_polynomial_checkpoint_gap_hiding.json}.

\section{Exact quantified target}

The universal necessity attacked in this run was:
\begin{quote}
For every uniform deterministic polynomial-time decider $A$ for SAT, there exists a
uniformly predeclared polynomial-size normalized schedule of internal checkpoints such
that the vector of mutable configurations at those checkpoints determines
$H(F,\alpha)$ on $\mathcal K_N$.
\end{quote}
Here ``mutable configuration'' excludes the immutable read-only input.  Including the
input would make ``exposure'' tautological and would still give no runtime lower bound.
The schedule is a uniformly computable list
\[
 C_N=(c_{N,1},\ldots,c_{N,m(N)})\subseteq\mathbb Q\cap[0,1],
 \quad m(N)\le N^d,
\]
whose rational descriptions and computation time are polynomial in $N$.  A cut $c$ of an
exact $Q(N)$-step run observes the configuration at step $\lfloor cQ(N)\rfloor$.

\section{Largest-gap wrapper theorem}

\begin{lemma}[interior grid cell]
Let $C\subseteq[0,1]$ contain at most $m$ points and put $M=8(m+1)$.  Some adjacent
points $a<b$ of $\{0,1\}\cup C$ satisfy $b-a\ge 1/(m+1)=8/M$.  Consequently there is an
integer $j$ with
\[
 a<\frac jM<\frac{j+1}{M}<b.
\]
\end{lemma}
\begin{proof}
The $|C|+1$ gaps have total length one, proving the first statement.  An interval of
length at least eight grid spacings contains a full grid cell after discarding boundary
cells; choosing any remaining cell gives both strict inequalities.
\end{proof}

\begin{theorem}[polynomial scheduled-cut hiding]
Let $A$ be any total uniform deterministic decider running in time at most $N^k$, and let
$C_N$ satisfy the conditions above.  There is a total uniform deterministic decider
$W_{A,C}$ such that, for every encoded input $x$ of length $N$:
\begin{enumerate}
\item $W_{A,C}(x)=A(x)$;
\item $W_{A,C}$ has an exact input-length-dependent halting time $Q(N)$ and runs in
polynomial time;
\item the vector of mutable configurations at every $c\in C_N$ is a function only of
$(N,A(x))$; and
\item its space is $O(s_A(N)+p_C(N)+\log N)$, where $s_A$ is the child's workspace and
$p_C$ is polynomial workspace sufficient to compute and sort the schedule.
\end{enumerate}
\end{theorem}

\begin{proof}
On a read-only input, first determine $N$ by a symbol-oblivious scan.  Compute and sort
$C_N$, choose the largest gap and the interior grid cell
$I=[j/M,(j+1)/M]$.  Choose an exactly clocked polynomial block length $B(N)$ dominating
the schedule computation, the $N^k$-step padded simulation of $A$, and a full cleanup of
all child work tapes.  Because the chosen cell is strictly interior, $j\ge1$, so the
first $j$ blocks suffice for the length-only schedule computation and source-independent
delay.  In block $j$, simulate $A$, retain its decision bit, erase all child-dependent
mutable state, and pad to exactly $B(N)$ steps.  Use source-independent delay for the
remaining $M-j-1$ blocks.  Halt after
\[
 Q(N)=M B(N)=8(m(N)+1)B(N)
\]
steps and output the retained bit.

No scheduled $c$ lies in $I$.  If $c<j/M$, then
$\lfloor cQ(N)\rfloor<jB(N)$ and the observed state is produced only by computation on
$N$ and $C_N$.  If $c>(j+1)/M$, the simulation and cleanup have finished, so the state
depends only on $N$, the cut index, and $A(x)$.  Exact clocks make these claims literal,
not asymptotic.  Computing the schedule, padding, simulation, and erasure are all charged;
$M$ and $B$ are polynomial.  The input is read-only and is not counted as exposed mutable
work state.
\end{proof}

\paragraph{Corollary.}
For the clean-core pair $F_0,F_1$ above, every scheduled state vector is identical because
the inputs have equal length and equal answer, while $H$ differs.  Therefore no function
of that vector alone determines $H$.  The target universal necessity is
\FALSIFIED.  The theorem does not say that an extractor also given $F$ cannot recompute
$H$; that stronger statement is precisely a complexity claim and is not obtained here.

\section{Reproducible finite audit and hostile controls}

The script
\path{experiments/review/polynomial_checkpoint_gap_hiding_audit.py} enumerated 288
schedules with $m\in\{1,2,3,4,8,16,32,64,128\}$, including even, clustered, endpoint,
and seeded-random schedules.  For each it verified a strict interior cell, absence of a
scheduled cut in the closed child interval, and equality of symbolic state fingerprints
for equal-length/equal-answer inputs.  It separately constructed and exhaustively checked
the clean-core pair.  Runtime was 2.259 seconds; no invocation exceeded 30 seconds.
\[
 \operatorname{SHA256}(\text{script})=
 \texttt{7f81303fe69489a2c83ba95e6c53f6f18100940f9a2b02f4dad5b21d5b01899b},
\]
\[
 \operatorname{SHA256}(\text{result})=
 \texttt{70a285cbbcb2248dfab06dbe4d1724f2a6f3e7b5394694a07af74db2c8f0c4f1}.
\]
The control ledger names 2SAT, Horn, affine/Tseitin, bounded treewidth, independent
products, padding, redundant clauses, alternate encodings, and wrapper-equivalent
programs.  The exact theorem uses only equal encoded length and equal decision bit, so
these controls cannot restore a scheduled-cut obligation.  This is a domain-independent
symbolic statement, not an empirical runtime comparison and not an asymptotic
extrapolation from the 288 cases.

\section{Relativization, algebrization, nonuniformity, and one-bit output}

The wrapper construction relativizes: an oracle call made by $A$ can be simulated inside
the hidden block.  That is evidence only for the negative theorem and cannot support a
P-versus-NP separation in view of Baker--Gill--Solovay.  No algebraic lower-bound method
is used; the algebrization barrier remains untouched.  A nonuniform checkpoint list of
polynomial description can likewise be hidden if the wrapper receives or contains that
list, so nonuniformity does not rescue this target.  Conversely, an exponentially large
or adaptively chosen list is outside the theorem.  The one-bit output is decisive: after
cleanup, a correct decider is obligated to retain only that bit.

\section{Post-quantum technique ledger}

\begin{longtable}{p{.30\linewidth}p{.62\linewidth}}
\toprule
Technique & Classification and consequence \\
\midrule
Schur squares; quadratic relations; Hilbert/syzygy invariants & \textbf{FAILED TRANSFER}
for this barrier.  These invariants concern a selected algebraic representation; the
wrapper theorem changes only timing and applies before any such representation is forced.
Hidden algebraic structure could instead be an algorithmic shortcut. \\
Lossiness/collision entropy & \textbf{ANALOGY/CANDIDATE}.  The scheduled observation map
has a large exact decision-fiber collision, but collision counting alone is not a time or
space lower bound. \\
Worst-case/average-case external extraction & \textbf{ANALOGY/CANDIDATE}.  Black-box
adaptive challenges might remove predeclaration, but no unconditional SAT transfer is
known and NP worst-case/average-case barriers must be audited. \\
Special soundness/collapse binding & \textbf{ANALOGY/CANDIDATE}.  Two conflicting
openings would need to yield a concrete forbidden witness in the same realization; no
such theorem is supplied. \\
Hash/random-oracle methods & \textbf{FAILED TRANSFER} for an unconditional proof unless
all hardness and random-oracle assumptions are removed. \\
\bottomrule
\end{longtable}

\section{CORPUS REVIEWED}

The repository tree was surveyed at Git commit
\texttt{2fc102d2d4ea1aceda0be99305a54490b5143435}.  Read directly: root
\path{main.tex}, \path{AGENTS.md}, \path{CODEX.md}, the complete PNP research skill and
mode/protocol references, every specification under \path{tasks/}, actual
\path{main4.tex} and \path{main5.tex}, newest verified MAIN
\path{research/main/2026-10-12_0850_MAIN_SYNTHESIS-0023.tex}, and materially relevant
historical MAIN material.  Specialist reports reviewed include HALTCANON1,
RESOURCEAXIS1, HEVALLB1, TDLG1, MAJORANTNOGO1, PRODSTREAM1, WRAPORBIT1, and
MODELBRIDGE-0026, plus the earlier LB-0006, MODELBRIDGE-0002,
CONSTRUCTION-STRENGTHEN9, and review STATUS artifacts surfaced by repository search.
The HEC, VSC-HEC, G-PPC-HEC, CRS-PPC, CCSPC, RCCN, and random-challenge exploration
branches, governance notes, and cited experiment/result paths were surveyed for a total
arbitrary-algorithm bridge.

The predecessor LSPV-5 check found no canonical artifact newer than SYNTHESIS-0023.
Because this repeated across more than two cycles, the synthesis automation prompt was
inspected and minimally repaired without changing its schedule or scientific target;
the evidence and rollback are in
\path{research/governance/2026-10-12_1414_LSPV5-WRITE-REPAIR-003.tex}.

\section{DUPLICATION / CONTRADICTION AUDIT}

Repository search for ``polynomially many checkpoints'', ``normalized-time checkpoints'',
and ``checkpoint vector exposure'' found no prior theorem hiding a whole polynomial
predeclared schedule.  MODELBRIDGE-0026 proves the one-fixed-cut case; RT-0009 covers a
fixed prefix; HALTCANON1 covers the final retained state.  CUTGAP1 strictly combines and
extends those negative results.  It does not contradict LB-0006's causal-cell route,
because that route requires an externally compiled/admissible checkpoint interface rather
than arbitrary internally normalized times.  It also does not contradict known
time--space tradeoffs, which impose resource hypotheses absent here.

\section{WHAT CHANGED RELATIVE TO ROOT MAIN}

The LSPV-1 specialty state is sharpened from ``one normalized cut can be pre-delayed'' to
``every uniformly predeclared polynomial-size normalized schedule can be avoided in one
largest gap.''  A valid observation-based route must now use an adaptive or exogenous
checkpoint, an all-times/peak property, or a model restriction that makes the charged
interval unreschedulable.  No executive claim, exponent, or proof status is strengthened
to a P-versus-NP result.

\section{ADOPTED / REJECTED / RETIRED PRIOR COMPONENTS}

\begin{itemize}
\item \textbf{ADOPTED:} read-only input convention, exact clocks, cleanup accounting,
same-realization discipline, and the equal-length/equal-answer hostile pair.
\item \textbf{ADOPTED:} MODELBRIDGE-0026's normalized-time attack as the one-cut base case.
\item \textbf{REJECTED:} inferring runtime or memory from the number of possible semantic
compatibility states.
\item \textbf{RETIRED:} every universal claim that some fixed polynomial predeclared set
of normalized cuts must expose $H$ in mutable state.
\item \textbf{RETAINED OPEN:} exogenous online challenges, adaptive cuts unknown to the
program, peak/all-time obligations, and a same-model transfer from any such property to an
unconditional lower bound.
\end{itemize}

\section{STALE OR OVERSTATED CLAIMS FOUND}

Any root or historical wording that treats a finite or polynomial observation schedule as
representative of an entire run is overstated unless the schedule is exogenous and
unreschedulable.  Likewise, a lower bound on a proof system, branching program, circuit,
communication protocol, or chosen representation is not a lower bound for every uniform
deterministic SAT decider without a total compiler preserving the charged resource in the
same realization.  The actual main4/main5 frameworks remain conditional and may not be
silently promoted.

\section{Highest-value missing implication}

The surviving target is an externally enforced interface theorem:
\begin{quote}
From every uniform deterministic polynomial-time decider for the NP-complete clean core,
construct in polynomial total cost an adaptive/exogenous observation or communication
experiment whose challenge cannot be predicted, delayed around, cleaned before, or
answered using only the decision bit, and whose success contradicts a model-matched
unconditional lower bound in the same realization.
\end{quote}
No such theorem is known here.  At least one load-bearing arrow is \OPEN; therefore P
versus NP remains open.

\begin{thebibliography}{9}
\bibitem{docker} J. D\"ocker, ``Monotone 3-Sat-(2,2) is NP-complete,'' 2019,
\url{https://arxiv.org/abs/1912.08032}.
\bibitem{fvm} L. Fortnow and D. van Melkebeek, ``Time-Space Tradeoffs for
Nondeterministic Computation,'' ECCC TR00-028, 2000,
\url{https://eccc.weizmann.ac.il/report/2000/028/}.
\bibitem{bgs} T. Baker, J. Gill, and R. Solovay, ``Relativizations of the
$\mathcal P=?\mathcal{NP}$ Question,'' \emph{SIAM J. Comput.} 4 (1975), 431--442,
\url{https://epubs.siam.org/doi/10.1137/0204037}.
\bibitem{aw} S. Aaronson and A. Wigderson, ``Algebrization: A New Barrier in Complexity
Theory,'' \emph{TOCT} 1(1), 2009,
\url{https://dl.acm.org/doi/10.1145/1490270.1490272}.
\end{thebibliography}

% END_OF_VERIFIED_REPORT::LSPV-1::CUTGAP1
\end{document}
